\chapter{Petunjuk dan Bentuk Penugasan}
\begin{warnbox}{Integritas Akademik}
\begin{center}
  Seluruh karya yang dikumpulkan harus merupakan hasil pemikiran dan kerja
  mahasiswa sendiri. Menyalin kode atau laporan dari sumber lain tanpa atribusi
  yang jelas adalah pelanggaran integritas akademik dan akan dikenai tindakan sesuai
  peraturan universitas.
  \end{center}
\end{warnbox}

\section{Petunjuk Tugas}
Tugas yang diberikan adalah menyangkut Model Relasional dan Query atas Model Relasional:
\begin{itemize}
    \itemsep0em
     \item Aljabar Relasional
     \item \textit{Structured Query Language (SQL})
\end{itemize}

Jawaban tugas dikemas dalam satu bab (bab 2), yang bagian-bagiannya telah disediakan dalam template ini. Isi Bab tersebut dapat dikembangkan sesuai kebutuhan, namun tidak untuk dikurangi.

\begin{deliverablebox}{Tugas terdiri atas empat bagian}
  \begin{enumerate}
    \itemsep0em
    \item menerjemahkan batasan dan proses bisnis ke dalam diagram ER (\textit{Entity Relationship Diagram}). Usulan disertai penjelasan proses pembentukan ERD
     \item merancang dan mengusulkan skema basis data relasional yang sesuai dengan ERD yang diusulkan
     \item Jika diberikan query, mengusulkan  representasi query dalam  ekspresi aljabar relasional
     \item Mengusulkan penyelesaian pertanyaan dalam bentuk query SQL. Bila ada lebih dari satu kemungkinan query (dengan operasi atau klausa yang berbeda), query alternatif juga perlu dituliskan
\end{enumerate}  
\end{deliverablebox}

\begin{infobox}{Ketentuan Tugas}
\begin{enumerate}
    \itemsep0em
     \item Pekerjaan dapat dikerjakan secara berkelompok dengan anggota kelompok maksimal 4 orang.
     \item Setiap anggota kelompok \textbf{harus} berkontribusi secara memadai, dan menguasai materi yang dihasilkan oleh kelompoknya.
     \item Pekerjaan harus disajikan secara \textit{jelas, lengkap dan rapi, termasuk penggunaan bahasa formal yang mudah dipahami}
     \item Pekerjaan harus dari pemikiran anggota tim, dan disepakati oleh semua anggota
\end{enumerate}  
\end{infobox}

\section{Detail Isi Penugasan}
\subsection{Domain Kasus}
Diinginkan untuk mengelola data akademik suatu perguruan tinggi dengan batasan sebagai berikut:
\begin{infobox}{Batasan Proses}
    \begin{itemize}
        \itemsep0em
     \item Setiap kegiatan akademik perlu dicatat dengan transparan dan mempertahankan kronologi/historis
     \item Desain penyimpanan data mengakomodasi hal-hal dan proses berikut:
     \begin{itemize}
         \itemsep0em \small
         \item Seorang mahasiswa diperkenankan mengulang matakuliah tanpa batas berapa kali
         \item Nilai matakuliah (apabila pernah mengulang) yang digunakan sebagai nilai untuk transkrip adalah \textbf{nilai pengambilan terakhir}
         \item Matakuliah ada yang bersifat \textsc{Wajib} dan ada \textsc{Pilihan}.
         \item Matakuliah dapat diprasyarati matakuliah lainnya
          \item Seorang dosen dapat mengajar suatu matakuliah di suatu semester lebih di satu kelas (paralel)
         \item Beberapa matakuliah dapat ditawarkan di semster genap ataupun semester ganjil
         \item Suatu kelas untuk suatu matakuliah diajar penuh oleh seorang dosen (tidak secara team teching)
         \item Nilai matakuliah dan bobot yang diberikan adalah adalah 
         \begin{center}
         \begin{multicols}{2}
            \begin{tabular}{c|c}
             Nilai & Bobot \\ \hline
           A   & 4,00 \\
            A-  & 3,75 \\
            A/B & 3,50 \\
            B+ & 3,25 \\
            B & 3,00 \\ 
            B-  & 2,75 \\
         \end{tabular}  
         \columnbreak
         \begin{tabular}{c|c}
             Nilai & Bobot \\ \hline
            B/C  & 2,50 \\
            C+ & 2,25 \\
            C & 2 \\
            C/D & 2,5 \\
            D & 1,00 \\ 
             D & 0,00 \\ 
         \end{tabular} 
         \end{multicols}
         \end{center}
         \item Rumus untuk menghitung IPK dari $n$ mata kuliah adalah:
         \[IPK = \frac{\sum_{i=1}^{n}{{\text{sks}}_i.{\text{bobot}}_i}}{\sum_{i=1}^{n}{{\text{sks}}_i}}\]
     \end{itemize}
    \end{itemize}
\end{infobox}

\subsection{Output/Outcome Penugasan}
Berikut adalah hal-hal minimal yang harus menjadi output/outcome penugasan.

\begin{deliverablebox}
    \begin{enumerate}
        \itemsep0em
     \item E-R diagram (ERD) yang sesuai hal-hal di atas! (dimulai dari identifikasi: entitas, relasi, kardinalitas, dsb.)
     \item Skema lengkap tabel-tabel basis data yang sesuai ERD yang dibuat; misal dalam bentuk \newline \texttt{Mhs(nim, nama, jenisKelamin, prodi,\dots) \newline Dosen(nip, \dots)}
     
     Dan tulis juga \textbf{perintah DDL SQL untuk definisi tabel-tabel} di atas (lengkap dengan semua \textit{constraint})!
     \item Perintah DML untuk pertanyaan-pertanyaan di bawah ini:
     \begin{enumerate}
         \itemsep0em \footnotesize
         \item Daftar nama mata kuliah pilihan dan mata kuliah yang memprasyaratinya (cukup info nama dan kode matakuliah saja).  Matakuliah yang tidak memiliki prasyarat tetap ditampilkan!
         \item Daftar Matakuliah pilihan yang diambil oleh mahasiswa Ilmu Komputer angkatan '2024' pada semester ganjil tahun akademik 2026/2027!
         \item Daftar Matakuliah pilihan yang favorit untuk mahasiswa Ilmu Komputer, yaitu tiga matakuliah dengan pengambil terbanyak dalam tiga tahun terakhir!
         \item Daftar mahasiswa yang memiliki Indeks Prestasi Semester ganjil tahun 2025/2026 tidak kurang dari 3,25! (pahami perbedaan antara IP Semester dan IP Kumulatif)
         \item Daftar mahasiswa dan IP Kumulatif yang diperolehnya saat ini! (matakuliah yang pernah diulang cukup dihitung sekali) 
     \end{enumerate}
     \item Query dalam bentuk ekspresi aljabar relasional untuk setiap pertanyaan/query yang diberikan!
    \end{enumerate}
\end{deliverablebox}

\vspace{.2cm}
Untuk setiap query di atas, langkah pengerjaan perlu diuraikan:
\begin{enumerate}
    \itemsep0em
     \item identifikasi tabel-tabel yang perlu dilibatkan
     \item penentuan kriteria join dan jenis join: \textsc{inner join, natural join, left/right join} atau \textsc{outer join}
     \item Perintah SQL yang valid. Apabila memerlukan \textsc{VIEW} sebagai temporary table, silahkan juga dituliskan secara runtut.
     \item apabila ada lebih dari satu SQL query yang \textit{equivalent}, akan menjadi nilai tambah apabila Anda bisa menuliskan serta query tersebut!
\end{enumerate}

\section{Contoh Soal dan Pengerjaannya}
Misal diberikan query:
\begin{tcolorbox}{}
  \begin{quote}
    \textit{Tampilkan daftar mata kuliah pilihan yang belum pernah diambil/diikuti oleh mahasiswa}
\end{quote}  
\end{tcolorbox}
maka contoh jawaban adalah sebagai berikut:
\begin{infobox}{Query nomor 1}
    \begin{enumerate}
        \itemsep0em \small
         \item \textbf{Tabel yang diperlukan dalam query:} \textsc{Matakuliah (Mk), PengambilanKuliah (Pk)}
         \item bila dengan operasi SET DIFFERENCE tidak ada proses join, sehingga tidak ada kriteria join. Sebagai alternatif, dapat digunakan LEFT JOIN, dengan kriteria \texttt{Mk.idMatakuliah= Pk.idMatakuliah} 
         \item \textbf{Query yang diperlukan} 
         \begin{lstlisting}[caption= pendekatan SET DIFFERENCE]
SELECT * FROM Matakuliah Mk WHERE (sifat='P') and idMatakuliah NOT IN 
    (SELECT idMatakuliah FROM PengambilanKuliah Pk)       
         \end{lstlisting}
         \item Alternatif SQL yang ekuivalen
         \begin{lstlisting}[caption=dengan LEFT JOIN]
SELECT Mk.* FROM Matakuliah Mk LEFT JOIN PesertaKuliah Pk ON (Mk.idMatakuliah=Pk.idMatakuliah) 
    WHERE Mk.sifat='P' and Pk.idMatakuliah is NULL
         \end{lstlisting}
      \begin{lstlisting}[caption=dengan EXISTS]
SELECT * FROM Matakuliah Mk WHERE (sifat='P') and 
    NOT EXISTS
    (SELECT idMatakuliah FROM PengambilanKuliah Pk 
        WHERE Pk.idMatakuliah=Mk.idMatakuliah)       
         \end{lstlisting}  
 \item Ekspresi Aljabar Relasional
 {\large
 \[
\begin{array}{c}
(\pi_{\text{idMatakuliah}}(\sigma_{\text{sifat='P'}}\left(\text{Matakuliah}) \right) - \\ \pi_{\text{idMatakuliah}}\left(\text{Matakuliah}\right))   \bowtie \text{Matakuliah}
\end{array}
 \]
}
{\large
 \[
 \begin{array}{c}
DIFF1 \leftarrow \pi_{\text{idMatakuliah}}(\sigma_{\text{sifat='P'}}\left(\text{Matakuliah}) \right)\\
 DIFF2 \leftarrow (DIFF1 - \pi_{\text{idMatakuliah}}\left(\text{Matakuliah}\right))  \\
DIFF2 \bowtie \text{Matakuliah}
 \end{array}
 \]
 }
    \end{enumerate}
\end{infobox}